Omdat Linux alleen de kernel is moet het gecombineerd worden met andere software om er een operatingsystem van te maken. Veel van die software komt van het GNU-project, maar dat is niet de enige leverancier. Er zijn vele open source projecten die ook software maken dat kan draaien op een GNU/Linux systeem. Door deze grote keuze vrijheid bestaan er natuurlijk verschillende meningen over wat wel en niet gewenst is binnen een operatingsystem. Wil je bijvoorbeel een besturingssysteem maken voor op een telefoon, op een desktop of op een server dan moeten er andere keuzes gemaakt worden. Los nog van het feit dat techneuten soms ook een eigen mening hebben en bepaalde stukken software beter vinden dan anderen. En zo ontstonden er verschillende verzamelingen van software met daarbij een Linux kernel.

En net als bij Unix werden die verzamelingen distributies genoemd. E\'en van de eerste distributies was Slackware, later kwamen SuSE, Debian, Red Hat, CentOS, Ubuntu en nog vele vele anderen. Een van de meest gebruikte \textquote{linux distributies} is waarschijnlijk Android dat op menig telefoon werkt.`

Het meest gebruikte Unix desktop systeem is ongetwijfeld Darwin het open source besturingssysteem van Apple dat de basisdd
vormt voor Mac OS X en op de telefoon is dat Android het systeem van Google dat een Linux kernel en andere open source
tools onderwater gebruikt.\par

En zo komen we na een lang verhaal aan het einde van deze simpele geschiedenis over Unix en Linux. Waarin hopelijk duidelijk is geworden waarom Linux Linux wordt genoemd, duidelijk is waarom Linux open source is en waarom delen zo'n belangrijk onderdeel is van de cultuur rond Unix en Linux systemen.

